\documentclass[xcolor={usenames,svgnames,x11names,dvipsnames,table}]{beamer}

\usetheme{SBUclass}

\usepackage{mypackages}

\title{Big Picture}
\author{Thomas Graf}
\institute{\texttt{lin220@thomasgraf.net}}
\date{LIN 220\\ Lecture 1}


\begin{document}
\unnumbered{
\begin{frame}
	\titlepage
\end{frame}
}

\begin{frame}{Two lessons}
    \begin{columns}
        \column{.45\linewidth}
        \centering
        \visible<2->{
        \begin{alertblock}{Lesson 1}
            Pogo sticks don't fly.
        \end{alertblock}
        \includegraphics[height=10em]{./img/pogostick}
        }
        \column{.45\linewidth}
        \visible<3->{
        \begin{alertblock}{Lesson 2}
            Wasps are strangely stupid.
        \end{alertblock}
        \includegraphics[height=10em]{./img/wasp}
        }
    \end{columns}
\end{frame}

\begin{frame}{An alternate history of flight}
    \begin{itemize}
        \item It is the 19th century, and airplanes aren't a thing yet.
        \item Three competitors at the first national flight competition:
    \end{itemize}

    \begin{tabular}{ccc}
        \begin{tikzpicture}
            \node (img) at (0,0) [visible on=<1->,
                                  label={[font=\bfseries]below:Icarus Inc}]
                {\includegraphics[width=8em]{./img/ikarus}};
            \draw[very thick, Red3, visible on=<2->] (img.south west) to (img.north east);
            \draw[very thick, Red3, visible on=<2->] (img.south east) to (img.north west);
        \end{tikzpicture}
        &
        \begin{tikzpicture}
            \node (img) at (0,0) [visible on=<3->,
                                  label={[visible on=<3->, font=\bfseries]below:Leonardo LLC}]
                {\includegraphics[width=8em]{./img/leonardo}};
            \draw[very thick, Red3, visible on=<4->] (img.south west) to (img.north east);
            \draw[very thick, Red3, visible on=<4->] (img.south east) to (img.north west);
        \end{tikzpicture}
        &
        \begin{tikzpicture}
            \node (img) at (0,0) [visible on=<5->,
                                  label={[visible on=<5->, font=\bfseries]below:Timmy}]
                {\includegraphics[width=3.5em]{./img/pogokid}};
        \end{tikzpicture}
    \end{tabular}
\end{frame}

\begin{frame}{An alternate history of flight [cont.]}
    \begin{columns}
        \column{.6\linewidth}
        \begin{itemize}
            \item After Timmy's victory,\\
                  \highlight{pogo sticks are all the rage}.
            \item Better and better pogo sticks\\
                  hit the market.
            \item By 1930, the US is the world's leading pogo stick nation.
        \end{itemize}
        \column{.5\linewidth}
        \includegraphics[width=10em]{./img/rising}

        \medskip
        \includegraphics[width=10em]{./img/fuckyeah}
    \end{columns}
\end{frame}

\begin{frame}{An alternate history of flight [cont.]}
    \begin{itemize}
        \item By 1952, the US is occupied by Japan and Nazi Germany. 
        \item Its pogo sticks were \highlight{no match for airplanes}.
        \item Nobody knows what happened to Timmy.
    \end{itemize}

    \begin{center}
    \includegraphics[width=24em]{./img/occupied_us}
    \end{center}
\end{frame}

\begin{frame}{Language technology = pogo sticks}
    \begin{itemize}
        \item Like flight, language technology is the future.
        \item Like pogo sticks, \highlight{hyped solutions} are not the final answer:
            \begin{itemize}
                \item Deep Learning
                \item Big Data
                \item embeddings
                \item RNNs, LSTMs
                \item seq2seq
            \end{itemize}
        \item Let's talk about them\ldots
    \end{itemize}
\end{frame}

\begin{frame}{The current hype: Deep Learning}
    \begin{itemize}
        \item One learning model is all over the media right now:\\
            \highlight{Deep Learning}
        \item Deep learning = very large and complex neural networks
        \item Neural networks imitate the human brain.
    \end{itemize}

    \pause
    \begin{block}{Standard model of the human brain}
        \begin{itemize}
            \item connected network of neurons
            \item input activates neurons, which start ``firing''\\
                (= emitting electrical current)
            \item current activates other neurons $\Rightarrow$ activation patterns
            \item learning = strengthening connection between specific neurons
        \end{itemize}
    \end{block}
\end{frame}

\begin{frame}{The perceptron}
    \begin{itemize}
        \item \textbf{input layer:} neurons that are sensitive to input
        \item \textbf{output layer:} neurons that represent output values
        \item \textbf{connections:} weighted links between input and output layer
        \item most activated output neuron represents decision
    \end{itemize}
    %
    \begin{center}
        \only<1>{
        \begin{tikzpicture}
            \node[minimum size=3em,draw,circle] (dear) at (0,0) {Dear};
            \node[minimum size=3em,draw,circle] (hi) [left=of dear] {Hi};
            \node[minimum size=3em,draw,circle] (emily) [right=of dear] {Emily};
            \node[minimum size=3em,draw,circle] (ham) [below=5em of hi] {ham};
            \node[minimum size=3em,draw,circle] (spam) [below=5em of emily] {spam};

            \foreach \Input/\Output/\Value in {%
                    dear/ham/0,
                    dear/spam/5,
                    hi/ham/3,
                    hi/spam/1,
                    emily/ham/10,
                    emily/spam/0}
                \draw[->] (\Input) to node [right, pos=.75] {\Value} (\Output);
        \end{tikzpicture}
        }%
        \only<2>{
        \begin{tikzpicture}
            \node[minimum size=3em,draw,circle,fill=blue!25] (dear) at (0,0) {Dear};
            \node[minimum size=3em,draw,circle,fill=blue!25] (hi) [left=of dear] {Hi};
            \node[minimum size=3em,draw,circle] (emily) [right=of dear] {Emily};
            \node[minimum size=3em,draw,circle,fill=red!10] (ham) [below=5em of hi] {ham};
            \node[minimum size=3em,draw,circle,fill=red!35] (spam) [below=5em of emily] {spam};

            \foreach \Input/\Output/\Value in {%
                    dear/ham/0,
                    dear/spam/5,
                    hi/ham/3,
                    hi/spam/1,
                    emily/ham/10,
                    emily/spam/0}
                \draw[->] (\Input) to node [right, pos=.75] {\Value} (\Output);
        \end{tikzpicture}
        }%
        \only<3>{
        \begin{tikzpicture}
            \node[minimum size=3em,draw,circle,fill=blue!25] (dear) at (0,0) {Dear};
            \node[minimum size=3em,draw,circle,fill=blue!25] (hi) [left=of dear] {Hi};
            \node[minimum size=3em,draw,circle,fill=blue!25] (emily) [right=of dear] {Emily};
            \node[minimum size=3em,draw,circle,fill=red!75] (ham) [below=5em of hi] {ham};
            \node[minimum size=3em,draw,circle,fill=red!35] (spam) [below=5em of emily] {spam};

            \foreach \Input/\Output/\Value in {%
                    dear/ham/0,
                    dear/spam/5,
                    hi/ham/3,
                    hi/spam/1,
                    emily/ham/10,
                    emily/spam/0}
                \draw[->] (\Input) to node [right, pos=.75] {\Value} (\Output);
        \end{tikzpicture}
        }
    \end{center}
\end{frame}

\begin{frame}{For the friendly neighborhood mathematician}
    \begin{itemize}
        \item Perceptrons are linear functions (matrix multiplication).
    \end{itemize}

    \begin{example}[Computing \emph{Hi Dear}]
        $$
        \begin{pmatrix}
            1 & 1 & 0\\
        \end{pmatrix}
        \otimes
        \begin{pmatrix}
            3 & 1\\
            0 & 5\\
            10 & 0\\
        \end{pmatrix}
        =
        \begin{pmatrix}
            3 & 6
        \end{pmatrix}
        $$
    \end{example}

    \begin{itemize}
        \item Since matrix multiplication is associative,\\
              every \highlight{multi-layer perceptron can be reduced to one layer}.
    \end{itemize}

    \visible<2->{%
    \begin{example}[Adding more weight to spam score]
        $$
        \begin{pmatrix}
            3 & 6\\
        \end{pmatrix}
        \otimes
        \begin{pmatrix}
            3 & 0\\
            0 & 1\\
        \end{pmatrix}
        =
        \begin{pmatrix}
            9 & 6\\
        \end{pmatrix}
        =
        \begin{pmatrix}
            1 & 1 & 0\\
        \end{pmatrix}
        \times
        \begin{pmatrix}
            9 & 1\\
            0 &  5\\
            30 & 0\\
        \end{pmatrix}
        $$
    \end{example}
    }
\end{frame}

\begin{frame}{Neural networks: The big picture}
    \begin{itemize}
        \item Modern neural networks are just the perceptron on steroids.
        \item There's a lot of jargon:
            \begin{itemize}
                \item RNNs
                \item LSTMs
                \item embeddings
                \item encoder/decoder
                \item seq2seq
            \end{itemize}
        \item Mathematically, modern neural networks intersperse linear functions (= perceptron layers) with non-linear functions.
        \item And \highlight{that's it}.\\
    \end{itemize}
\end{frame}

\begin{frame}{Neural networks aren't a magic bullet}
    Neural networks are not perfect:
    \begin{itemize}
        \item \textbf{Data hungry}\\
              If you don't have tons of data, don't even try.
        \item \textbf{Resource hungry}\\
              Large networks take forever to train.
        \item \textbf{Black boxes}\\
              Nobody knows what they do.
              It's \highlight{trial and error}.
        \item \textbf{Do not scale}\\
              If your objectives change, you're back to square~1.
        \item \textbf{No safeties}\\
              If something goes wrong, it often goes really wrong.
    \end{itemize}
\end{frame}

\begin{frame}{Neural networks in one picture}
    \centering
    \includegraphics[width=18em]{./img/xkcd}
\end{frame}

\begin{frame}{A common reply}
    \begin{quotation}
        We're not aiming for perfection,\\
        it just has to be good enough.
    \end{quotation}
    \hspace{14em} Every engineer ever

    \bigskip
    \textbf{My reply}
    \begin{enumerate}
        \item It's still a bad choice for engineering:
            \begin{itemize}
                \item expensive (resources, time, labor)
                \item scales badly
                \item not modular
            \end{itemize}
        \item Your notions of ``good enough'' are too limited:
            \begin{itemize}
                \item precision
                \item recall
                \item F-score
            \end{itemize}
              They all ignore \highlight{error quality}.
    \end{enumerate}
\end{frame}

\begin{frame}{Let's go to wasp school}
    \begin{columns}
        \column{.45\linewidth}
        \begin{alertblock}{Lesson 2}
            Wasps are strangely stupid.
        \end{alertblock}
        \column{.45\linewidth}
        \includegraphics[height=10em]{./img/wasp}
    \end{columns}
    \begin{center}
        \visible<2>{
            \LARGE\bfseries \color{SteelBlue4}Volunteer needed!
        }
    \end{center}
\end{frame}

\begin{frame}{What we learned from wasp school}
    \begin{itemize}
        \item Users endow systems with human-like qualities.
        \item When human biases are violated, the illusion breaks down.
        \item Breaking the illusion is jarring.
    \end{itemize}

    \begin{block}{The true task of language technology}
        \begin{itemize}
            \item Trick humans into considering you human-like.
            \item Minimize errors that violate human biases.
        \end{itemize}
    \end{block}
\end{frame}

\begin{frame}{Biases in human cognition: a quick experiment}
    \begin{center}
        \small
        \begin{tabular}{cccccccc}
            \begin{tikzpicture}
                \fill[blue] (0,0) circle (1em);
            \end{tikzpicture}
            &
            \begin{tikzpicture}
                \fill[purple] (0,0) circle (1em);
            \end{tikzpicture}
            &
            \begin{tikzpicture}
                \fill[blue] (0,0) rectangle (2em,2em);
            \end{tikzpicture}
            &
            \begin{tikzpicture}
                \fill[purple] (0,0) rectangle (2em,2em);
            \end{tikzpicture}
            &
            \begin{tikzpicture}
                \fill[brown] (0,0) circle (1em);
            \end{tikzpicture}
            &
            \begin{tikzpicture}
                \fill[brown] (0,0) rectangle (2em,2em);
            \end{tikzpicture}
            &
            \begin{tikzpicture}
                \fill[brown] (0em,0em) -- (2em,0em) -- (1em,2em) -- cycle;
            \end{tikzpicture}
            \\[6pt]
            \emph{blip}
            &
            not \emph{blip}
            &
            \emph{blip}
            &
            not \emph{blip}
            &
            \visible<2->{%
            \emph{blip}
            }
            &
            \visible<3->{%
            \emph{blip}
            }
            &
            \visible<4->{%
            \emph{blip}
            }
            \\[6pt]
            not \emph{gnok}
            &
            not \emph{gnok}
            &
            \emph{gnok}
            &
            \emph{gnok}
            &
            \visible<6->{%
            \emph{gnok}
            }
            &
            \visible<7->{%
            \emph{gnok}
            }
            &
            \visible<8->{%
            \emph{gnok}
            }
            \\[6pt]
            not \emph{bnik}
            &
            not \emph{bnik}
            &
            \emph{bnik}
            &
            not \emph{bnik}
            &
            \visible<10->{%
            \emph{bnik}
            }
            &
            \visible<11->{%
            \emph{bnik}
            }
            &
            \visible<12->{%
            \emph{bnik}
            }
            \\[6pt]
            \emph{glop}
            &
            \emph{glop}
            &
            not \emph{glop}
            &
            \emph{glop}
            &
            \visible<14->{%
            \emph{glop}
            }
            &
            \visible<15->{%
            \emph{glop}
            }
            &
            \visible<16->{%
            \emph{glop}
            }
            \\[6pt]
            \emph{blok}
            &
            \emph{blok}
            &
            \emph{blok}
            &
            \emph{blok}
            &
            \visible<18->{%
            not \emph{blok}
            }
            &
            \visible<19->{%
            not \emph{blok}
            }
            &
            \visible<20->{%
            not \emph{blok}
            }
        \end{tabular}
    \end{center}

    \begin{description}
        \item[blip] {\visible<5->{not red}}
        \item[gnok] {\visible<9->{brown or rectangular}}
        \item[bnik] {\visible<13->{blip and gnok}}
        \item[glop] {\visible<17->{if bnik and not brown, then not rectangular}}
        \item[blok] {\visible<21->{bnik or glop, but not both}}
    \end{description}
\end{frame}

\begin{frame}{Human language bias and unreasonable expectations}
    Human bias is a much bigger issue for language than for, say, cars.
    %
    \begin{center}
        \begin{tabular}{cccc}
            \includegraphics[height=8em]{./img/data} &
            \includegraphics[height=8em]{./img/c3po} &
            \includegraphics[height=8em]{./img/hal9000} &
            \includegraphics[height=8em]{./img/terminator} \\
            \includegraphics[height=8em]{./img/robby} &
            \includegraphics[height=8em]{./img/tachikoma} &
            \includegraphics[height=8em]{./img/brainiac} &
            \includegraphics[height=8em]{./img/shodan}
        \end{tabular}
    \end{center}
\end{frame}

\begin{frame}{Robots' Narrow Range of Language Competence}
    \begin{tikzpicture}
        \begin{scope}[visible on=<2->]
            \node (perfect) at (0,0) {%
                {
                    \includegraphics[width=5em]{img/wheatley}}%
                };
            \node (perfect-text) [right=2em of perfect] {completely human};
        \end{scope}

        \begin{scope}[visible on=<3->]
            \node (good) [below=1em of perfect] {%
                {
                    \includegraphics[width=5em]{img/robby_cropped}}%
                };
            \node (good-text) [right=2em of good] {perfect but weird voice};
        \end{scope}

        \begin{scope}[visible on=<4->]
            \node (arnold) [below=1em of good] {%
                {
                    \includegraphics[width=5em]{img/terminator}}%
                };
            \node (arnold-text) [right=2em of arnold] {Arnold};
        \end{scope}

        \draw[->,red,very thick] (perfect.north west) -- (arnold.south west);
    \end{tikzpicture}
\end{frame}

\begin{frame}{Fake it till you make it!}
    \begin{itemize}
        \item Language technology is still largely smokes and mirrors.
        \item User-facing software has to fool the user.
        \item Current methods can't do that, they will trip up in weird ways.
    \end{itemize}

    \begin{block}{Be wrong, not weird}
        \begin{itemize}
            \item 90\% performance can be better than 99\%.
            \item It depends on how unnatural the errors are.
        \end{itemize}
    \end{block}
\end{frame}

\begin{frame}{Wrapping up}
    \begin{itemize}
        \item Don't get me wrong, NLP has achieved amazing things.
        \item Neural networks have revolutionized the field.
        \item But NNs are no magic fix.
        \item For long-term success,\\
              \highlight{language technology needs language science}.  \end{itemize}

    \begin{columns}
        \column{.45\linewidth}
        \centering
        \visible<2->{
        \begin{alertblock}{Lesson 1}
            Pogo sticks don't fly.
        \end{alertblock}
        }
        \column{.45\linewidth}
        \visible<3->{
        \begin{alertblock}{Lesson 2}
            Wasps are strangely stupid.
        \end{alertblock}
        }
    \end{columns}
\end{frame}

\end{document}
